% chapters/08/section-04.tex
\section{Kapan Mengejar dan Kapan Berhenti}
\label{sec:berhenti}

Perilaku mengejar memerlukan dua batas, bukan satu. Bagian ini menjelaskan
keduanya dan persoalan yang muncul apabila batas yang kedua dilupakan.

\subsection{Jarak Henti dan Getaran di Atas Pemain}

Apabila NPC terus mengejar sampai jaraknya nol, kecepatannya selalu mengarah
ke pemain, dan pada setiap bingkai ia berpindah sejauh kecepatan dikalikan
selang waktu. Perpindahan itu tidak dapat berhenti tepat pada posisi pemain:
NPC akan melampauinya, lalu berbalik, melampauinya lagi, dan seterusnya.

Pada kecepatan 150 piksel per detik dan laju 60 bingkai per detik, perpindahan
setiap bingkainya 2,5 piksel. NPC akan bergetar di sekitar pemain dengan
simpangan sekitar 2,5 piksel. Getaran itu tidak hanya jelek dipandang; ia juga
membuat NPC tampak menempel pada pemain dan mengganggu pemeriksaan tabrakan.

Penanganannya adalah menentukan jarak henti: apabila jarak ke pemain sudah
kurang dari nilai itu, NPC berhenti.

\begin{kodekeluaran}
if self.distance <= self.stop_distance:
    self.velocity = pygame.Vector2(0, 0)
elif self.distance <= self.radius:
    self.velocity = to_player.normalize() * self.speed
else:
    self.velocity = pygame.Vector2(0, 0)
\end{kodekeluaran}

\begin{tipsbox}
Perhatikan hubungan antara jarak henti dan perpindahan setiap bingkai. Agar
getaran tidak terjadi, \textbf{jarak henti harus lebih besar daripada
perpindahan satu bingkai}, yaitu kecepatan dibagi laju bingkai. Pada contoh di
atas, perpindahan setiap bingkainya 2,5 piksel sedangkan jarak hentinya 30
piksel, sehingga aman. Akan tetapi, apabila kecepatannya dinaikkan menjadi
3.000 piksel per detik, perpindahan setiap bingkainya menjadi 50 piksel dan
melampaui jarak henti, sehingga NPC akan melewati daerah henti tanpa pernah
berhenti di dalamnya.
\end{tipsbox}

\subsection{Dua Batas, Empat Keadaan}

Dengan adanya kedua batas itu, keputusan NPC memiliki empat kemungkinan.
Tabel~\ref{tab:empat-keputusan} memuatnya.

\begin{table}[htbp]
\centering
\caption{Empat kemungkinan keputusan NPC menurut jaraknya terhadap pemain}
\label{tab:empat-keputusan}
{\setlength{\tabcolsep}{4pt}\footnotesize%
\begin{tabularx}{\textwidth}{@{}L{4.4cm} L{2.8cm} Y@{}}
\toprule
\headrow \textbf{Jarak ke pemain} & \textbf{Keputusan} & \textbf{Kecepatan} \\
\midrule
Lebih besar daripada jari-jari persepsi & Berdiam &
Nol. Pemain tidak terlihat. \\
\addlinespace
Di antara jari-jari dan jarak henti & Mengejar &
Arah menuju pemain, besarnya tetap. \\
\addlinespace
Lebih kecil daripada jarak henti & Berhenti &
Nol. Pemain sudah tercapai. \\
\bottomrule
\end{tabularx}}
\tabsource{Baris kedua adalah keadaan yang menghasilkan perilaku mengejar,
yaitu perilaku yang menjadi pokok bab ini.}
\end{table}

\subsection{Penghalang pada Batas Persepsi}

Persoalan kecil muncul pada batas jari-jari persepsi. Apabila pemain berada
tepat pada batas itu---misalnya berjalan mengelilingi NPC---keputusan NPC
berganti-ganti antara mengejar dan berdiam pada setiap bingkai, karena jaraknya
naik turun di sekitar nilai batas. Gejalanya terlihat sebagai NPC yang bergetar
dan berkedip warnanya.

Penanganannya sederhana dan telah lama dipakai pada perkakas lain: pakai
\textbf{dua nilai batas}, bukan satu. NPC mulai mengejar ketika jaraknya
kurang dari jari-jari persepsi, tetapi tidak berhenti mengejar sebelum jaraknya
melampaui nilai yang sedikit lebih besar.

\begin{kodekeluaran}
if self.chasing:
    batas = self.radius * 0.85      # bertahan sedikit lebih lama
else:
    batas = self.radius

if self.distance > batas:
    self.chasing = False
elif self.distance > self.stop_distance:
    self.chasing = True
\end{kodekeluaran}

Cara itu disebut \istilah{hysteresis}, dan dipakai pada banyak alat pengatur
untuk mencegah pergantian keadaan yang berlebihan.

\begin{catatabox}
Perhatikan variabel \code{self.chasing} pada potongan kode di atas. Variabel itu
mengingat keputusan pada bingkai sebelumnya, sehingga keputusan pada bingkai
ini tidak hanya bergantung pada jarak, melainkan juga pada keadaan sebelumnya.
Itulah keadaan pertama yang dimiliki sebuah NPC. Menambah keadaan lain, misalnya
berpatroli dan menyerang, beserta aturan perpindahan di antaranya, akan
menghasilkan mesin keadaan yang dibahas pada Bab~IX.
\end{catatabox}
